\chapter{Tables and data grids}
\label{ch:tables}

\begin{aside}
A table is a 2-D structure with columns and rows. The data is
homogeneous within columns. The layout is regular. The
challenge is keeping a 1000-row table interactive.
\end{aside}

\section{Column model}

\begin{lstlisting}
struct Column {
    std::string title;
    int width;          // or 0 = flex
    std::function<Element(const Row&)> render;
};

std::vector<Column> cols = {
    {"Name", 20, [](auto& r) { return text(r.name); }},
    {"Age",   5, [](auto& r) { return text(std::to_string(r.age)); }},
    {"Email", 0, [](auto& r) { return text(r.email); }},
};
\end{lstlisting}

\section{Header row}

A frozen header row sits above the scrollable body:

\begin{lstlisting}
auto table = column(
    table_header(cols),
    virtual_list(rows, [&](int i) {
        return table_row(cols, rows[i]);
    })
);
\end{lstlisting}

\section{Selection}

Single or multi-row selection:

\begin{lstlisting}
auto selected = signal<std::set<int>>({});
\end{lstlisting}

Click toggles; Shift-click ranges; Ctrl-click multi. The
visual cue: highlight the selected row's background.

\section{Sorting}

Click a header to sort by that column. The sort key is a
signal:

\begin{lstlisting}
auto sort_by = signal<int>(0);
auto sort_dir = signal<bool>(true);  // ascending

auto sorted = memo([&] {
    auto v = items;
    std::sort(v.begin(), v.end(), ...);
    return v;
});
\end{lstlisting}

The memo recomputes only when sort signals change.

\section{Filtering}

A search box drives a filter:

\begin{lstlisting}
auto filter = signal<std::string>("");
auto filtered = memo([&] {
    auto q = ctx.get(filter);
    std::vector<Row> out;
    for (auto& r : items)
        if (matches(r, q)) out.push_back(r);
    return out;
});
\end{lstlisting}

The sort and filter compose: filter first (smaller list),
sort second.

\section{Column resizing}

A drag handle on the column boundary lets the user resize.
Track the drag with mouse events:

\begin{lstlisting}
on_drag([&](int dx) {
    auto w = ctx.get(col_widths)[i];
    ctx.set(col_widths, /* updated vector */);
});
\end{lstlisting}

\section{Frozen columns}

A frozen leftmost column stays visible during horizontal
scroll. Implement as two side-by-side virtual lists sharing
a scroll signal for vertical position.

\section{Cell editing}

Click a cell to edit in place. The cell becomes a text
input; commit on blur or enter; cancel on escape.

\begin{lstlisting}
auto cell = ctx.get(editing) == (r, c)
    ? text_input(value)
    : text(value) | on_click([&] {
        ctx.set(editing, {r, c});
      });
\end{lstlisting}

\section{Recap}

\begin{recap}
\begin{itemize}[leftmargin=*]
  \item Columns: title, width, render function.
  \item Header above, virtualised body below.
  \item Selection, sort, filter as signals.
  \item Column resize via drag.
  \item In-place cell editing.
\end{itemize}
\end{recap}

\section*{Workshop}

\begin{enumerate}[leftmargin=*]
  \item Build a 10-column 10,000-row table.
  \item Add sort by clicking headers.
  \item Implement in-place editing for one column.
\end{enumerate}
